\section{From excess stability to a label}\label{sec:extraction}

This section contains the analytic statement used in the reduction. Its
constants do not depend on the dimension of the affine label space.
The functions may depend on an additional random parameter $\omega$, whose
law is independent of all rows, gate tables, and the random slice. In the application,
$\omega$ will be a tuple of right endpoints of the game.

Let $\phi_\omega:\bits^{\mathcal A}\to[-1,1]$ be odd, meaning
$\phi_\omega(-w)=-\phi_\omega(w)$, and set
\[
 f_\omega(R,F)=\phi_\omega(W(R,F)),\qquad
 \bar f_\omega=T_r f_\omega,\qquad r=1-\mu/d.
\]
The word, tree, and random slice $\mathcal D$ are those of
\Cref{sec:tree}. Write $F_{\mathrm{out}}$ for the tables outside
$\mathcal D$. Once $\omega,\mathcal D,R,F_{\mathrm{out}}$ are fixed,
$\bar f_\omega$ is a function of the independent slice-table bits.

For a node $v\in\mathcal D$, use the free shift coordinates
$(L_{gb})_{g\in[s],b\in\{0,1\}}$ and the remaining context
$\Gamma\in H^{k-2s}$ from \Cref{lem:shift-coordinates}. A hat marked
``shifts'' expands only these $2s$ free rows. For a label $q\in\mathcal A$
(equivalently, $q\in H^*$ with $q(\one)=1$),
define
\begin{equation}\label{eq:mass}
 m_q(f;\Gamma,F_{\mathrm{out}})
 =\E_{F_{\mathcal D}}
 \sum_{\substack{Q\in(H^*)^{2s}\\
                  Q_{g0}+Q_{g1}=q\ (g\in[s])}}
 \abs{\widehat f^{\,\mathrm{shifts}}(Q)}^2.
\end{equation}
The dependence on $v,\mathcal D$, and $\omega$ is suppressed. The sets of
indices in this definition are disjoint as $q$ varies. For a fixed $q$,
the mass includes every way to distribute that common sum between the
two children of each pair. We seek such mass for the original smoothed
function $\bar f_\omega$. The unsplit projection and a table-degree
cutoff will serve only to find a label carrying this mass; neither
projection is part of the conclusion. Keeping all the indices with a
given common sum will let us compare masses under affine pullback in
\Cref{sec:decoding}.

In particular, if
$|f|\le1$, then at every fixed context
\begin{equation}\label{eq:mass-bound}
 0\le\sum_{q\in\mathcal A}m_q(f;\Gamma,F_{\mathrm{out}})\le1.
\end{equation}

\begin{proposition}[Dimension-independent extraction]\label{prop:extraction}
Fix $t\in(0,1)$, $0<\delta<1-B(t)$, and $0<\mu<1/2$.
There are tree parameters $s\ge2,m\ge1$ and constants $p_*,\theta>0$,
depending only on $t,\delta,\mu$, with the following property.
For any affine label space $\mathcal A$ and any family $\phi_\omega$ as above,
if
\begin{equation}\label{eq:excess-stability}
 \E_{\omega,\mathcal D,R,F_{\mathrm{out}}}
 \Stab_t^{F_{\mathcal D}}(\bar f_\omega)
 \ge B(t)+\delta,
\end{equation}
then
\begin{equation}\label{eq:extraction-event}
 \Pr\left[
   \max_{q\in\mathcal A}
   m_q(\bar f_\omega;\Gamma,F_{\mathrm{out}})\ge\theta
 \right]\ge p_*.
\end{equation}
Here we sample $\omega$, the random slice $\mathcal D$, and a uniform
node $v\in\mathcal D$; conditional on these, we sample
$F_{\mathrm{out}}$ and $\Gamma$ independently and uniformly.
\end{proposition}

\begin{proof}
Excess stability first produces an influential slice entry. We retain
that influence after the unsplit projection, discard the part even in
its gate table, and express the remaining influence as an affine hash
of vectors indexed by labels. The hashing lemma then forces one label
to carry substantial conditional mass.

All constants chosen below are independent of $\mathcal A$ and of the
law of $\omega$. The degree cutoff $K$ and influence threshold $\tau$
are chosen first. We leave $s,m$ unspecified during the estimates:
Step~4 will choose $s$ to control the even and hashed contributions, and
then $m$ to control the cost of the unsplit projection.

\paragraph{Step 1: find an influential slice entry.}
Simultaneously negating the slice tables negates the word. Since
$\phi_\omega$ is odd and smoothing preserves this identity,
$\bar f_\omega$ has mean zero in the slice bits for every fixed
$\omega,\mathcal D,R,F_{\mathrm{out}}$. It also lies in $[-1,1]$.
Apply \Cref{lem:mis} with $\xi=\delta/2$, obtaining $K\ge1$ and $\tau>0$.
We may decrease $\tau$ to ensure $\tau\le1$.
Since stability is at most one, \eqref{eq:excess-stability} implies that,
with probability at least
\[
 \eta=\delta/2,
\]
some entry $z\in\F^s$ of some $v\in\mathcal D$ has
$\Inf_{v,z}^{\le K}(\bar f_\omega)\ge\tau$.
Indeed, if the probability of this event is $p$, the expected stability is
at most $B(t)+\delta/2+p$.

Let $S_v$ be the orthogonal row projection onto indices unsplit at $v$.
Put
\[
 H_v=P_{\le K}^{F_{\mathcal D}}S_v\bar f_\omega.
\]
By \Cref{lem:unsplit}, the required bound on the projection loss is
\begin{equation}\label{eq:choose-m}
 \E\sum_{v\in\mathcal D}
 \norm{(1-S_v)\bar f_\omega}_{2,R,F}^2
 \le\frac{s}{2em\mu}\le\frac{\eta\tau}{8}.
\end{equation}
For any fixed $s$, the last inequality can be ensured by increasing $m$.
The same expectation is the expectation over
$\omega,\mathcal D,R,F_{\mathrm{out}}$ of the sum of the corresponding
conditional squared norms in $F_{\mathcal D}$. Markov's inequality shows
that this sum is at most $\tau/4$ except on an event of probability
$\eta/2$. On the intersection with the influential-entry event, the
triangle inequality for the projection defining that influence gives
\[
 \sqrt{\Inf_{v,z}(H_v)}
 \ge\sqrt{\Inf_{v,z}^{\le K}(\bar f_\omega)}
      -\norm{(1-S_v)\bar f_\omega}_{2,F_{\mathcal D}}
 \ge\frac{\sqrt\tau}{2}.
\]
Consequently,
\begin{equation}\label{eq:large-max-sum}
 \E\sum_{v\in\mathcal D}\max_z\Inf_{v,z}(H_v)
 \ge\frac{\eta\tau}{8}=:\Lambda.
\end{equation}

We will also use
\begin{equation}\label{eq:total-influence}
 \E\sum_{v\in\mathcal D}\sum_z\Inf_{v,z}(H_v)\le K.
\end{equation}
To verify it, integrate over all rows. Each $S_v$ is an orthogonal
projection in the rows and commutes with the slice-bit Fourier projections.
It can therefore only decrease the integrated influence at its entries.
After dropping each $S_v$ separately, the left side is bounded by the
total influence of $P_{\le K}^{F_{\mathcal D}}\bar f_\omega$, averaged over
the remaining variables. Equation~\eqref{eq:degree-influence} and
$|\bar f_\omega|\le1$ give \eqref{eq:total-influence}.
No pointwise boundedness of $H_v$ is asserted or needed.

\paragraph{Step 2: discard even dependence on a gate.}
Fix $v\in\mathcal D$, and split $H_v$ into its even and odd parts under
negation of $F_v$ alone. These parts have disjoint Fourier supports in the
slice bits, also after retaining the monomials containing any specified
entry. Thus
\[
 \Inf_{v,z}(H_v)=
 \Inf_{v,z}(H_v^{\mathrm{even}})
 +\Inf_{v,z}(H_v^{\mathrm{odd}}).
\]
By the sign action in \Cref{lem:row-symmetry}, the even part has
$q_v(\one)=0$ on row support. Since the index is unsplit, all its child
indices are either $0$ or $q_v$, so they too take value zero on $\one$.
The domain-translation action in that lemma shows that the even part is
invariant under translations of the domain of $F_v$, with the rows and
outside tables fixed. These translations act transitively on its entries.
All entry influences at $v$ are therefore equal. By
\eqref{eq:total-influence},
\begin{equation}\label{eq:even-influence}
 \E\sum_{v\in\mathcal D}
 \max_z\Inf_{v,z}(H_v^{\mathrm{even}})\le\frac{K}{2^s}.
\end{equation}

\paragraph{Step 3: express odd influences as affine hashes.}
Fix $\omega,\mathcal D,v,\Gamma,F_{\mathrm{out}}$. The remaining row
coordinates $L_{gb}$ are independent uniform elements of $H$. On the odd
part, the common pair sum $q$ is a label, by
\eqref{eq:labels-dual}, and hence is nonzero. An unsplit shift index with
common sum $q$ puts $q$
in the second child on a subset $I\subseteq[s]$ and in the first child
elsewhere. Let $a_{q,I}(F_{\mathcal D})$ be its coefficient in
$H_v^{\mathrm{odd}}$.

Set
\[
 L_\Sigma=\sum_{g=1}^sL_{g0},\qquad J_g=L_{g0}+L_{g1},\qquad
 \mathbf J(q)=(q(J_1),\ldots,q(J_s)).
\]
The $s+1$ forms $L_\Sigma,J_1,\ldots,J_s$ are independent and uniform: the
displayed linear map from $H^{2s}$ onto $H^{s+1}$ is surjective. The
remaining $s-1$ coordinates do not occur in the unsplit expansion.
The character of the shift index $(q,I)$ is
\[
 \chi_q(L_\Sigma)(-1)^{\sum_{g\in I}q(J_g)}.
\]
On the other hand, the domain-translation identity gives
\[
 a_{q,I}(F_v(\cdot+z),F_{\mathcal D\setminus\{v\}})
 =(-1)^{\sum_{g\in I}z_g}a_{q,I}(F_{\mathcal D}).
\]
With $\Psi_q=\sum_{I\subseteq[s]}a_{q,I}$ we obtain
\begin{equation}\label{eq:odd-hash-form}
 H_v^{\mathrm{odd}}
 =\sum_{q\in\mathcal A}\chi_q(L_\Sigma)
 \Psi_q\bigl(F_v(\cdot+\mathbf J(q)),
                   F_{\mathcal D\setminus\{v\}}\bigr).
\end{equation}
The functions $a_{q,I}$ for distinct $I$ are orthogonal in the uniform
slice bits: a domain translation whose two characters differ preserves
their inner product and changes its sign. Consequently,
\begin{equation}\label{eq:psi-mass}
 \norm{\Psi_q}_{2,F_{\mathcal D}}^2
 =\sum_I\norm{a_{q,I}}_{2,F_{\mathcal D}}^2
 \le m_q(\bar f_\omega;\Gamma,F_{\mathrm{out}}).
\end{equation}
The inequality follows because the unsplit row projection is determined
by the free shift index, and the degree and parity projections are
orthogonal projections in the slice bits at each fixed context.
Every $\Psi_q$ has degree at most $K$.

Equation~\eqref{eq:odd-hash-form} has converted the free-shift characters
into translations of a gate table. We next collect its Boolean Fourier
coefficients into vectors whose squared norms are entry influences.
These Fourier coefficients are indexed by sets of table entries; the
earlier index $I\subseteq[s]$ instead specified a character of the
table's domain $\F^s$.

Denote by
$\widehat\Psi_q(B,D)$ the slice Fourier coefficient with entry set
$B\subseteq\F^s$ at gate $v$ and entry set $D$ on the other slice gates.
For $z\in\F^s$ define a real vector $c(q,z)$, with coordinates indexed by
pairs $(A,D)$ where $0\in A\subseteq\F^s$, by
\[
 c(q,z)_{A,D}=\widehat\Psi_q(z+A,D).
\]
Using entry sets relative to $z$ makes this vector space the same for
every $z$ and every label $q$.
Set $c(q,z)=0$ for functionals $q$ which are not labels. For the influence
at $z$, a slice Fourier set $B$ containing $z$ has the unique form
$B=z+A$ with $0\in A$. Translation by $\mathbf J(q)$ sends its coefficient
to $\widehat\Psi_q(z+A+\mathbf J(q),D)$. Hence, exactly,
\begin{equation}\label{eq:influence-hash}
 \Inf_{v,z}(H_v^{\mathrm{odd}})
 =\norm*{\sum_q\chi_q(L_\Sigma)c(q,z+\mathbf J(q))}^2.
\end{equation}
Writing $w_q=\sum_z\norm{c(q,z)}^2$ and $M_v=\sum_qw_q$, we have
\begin{equation}\label{eq:wq-mq}
 w_q=\sum_{B,D}|B|\abs{\widehat\Psi_q(B,D)}^2
 \le K\norm{\Psi_q}_2^2
 \le K m_q(\bar f_\omega;\Gamma,F_{\mathrm{out}}).
\end{equation}
The factor $|B|$ counts the entries whose influences contain that
Fourier monomial; it is at most the monomial's total degree, which is
at most $K$.
Thus $M_v\le K$ at every context. The total-mass identity in
\Cref{lem:hashing}, together with \eqref{eq:influence-hash}, also gives
\begin{equation}\label{eq:sum-M}
 \E\sum_{v\in\mathcal D}M_v
 =\E\sum_{v\in\mathcal D}\sum_z
     \Inf_{v,z}(H_v^{\mathrm{odd}})\le K.
\end{equation}
For this identity the free coordinates are integrated separately for each
node; their different coordinate systems do not affect the uniform row law.

\paragraph{Step 4: force a heavy label.}
Choose $\zeta>0$ with $\zeta K\le\Lambda/4$. Apply
\Cref{lem:hashing}, obtaining $\gamma>0$ and a lower bound on $s$.
Choose $s\ge2$ large enough for that lemma and for
$K2^{-s}\le\Lambda/4$. Then choose $m$ as in \eqref{eq:choose-m}.
Finally put
\begin{equation}\label{eq:extraction-constants}
 M_0=\frac{\Lambda}{4\cdot2^m},\qquad
 \theta=\frac{\gamma M_0}{K},\qquad
 p_* =\frac{\Lambda}{4K\,2^m}.
\end{equation}

For each node, condition on $\omega,\mathcal D,v,\Gamma,F_{\mathrm{out}}$
and integrate the free shifts. Partition these contexts into three classes:
\begin{enumerate}[label=(\roman*)]
\item $w_q\le\gamma M_v$ for every $q$;
\item the first condition fails and $M_v<M_0$;
\item $M_v\ge M_0$ and $w_q>\gamma M_v$ for some $q$.
\end{enumerate}
On class (i), the hashing lemma bounds the expected maximum odd influence
by $\zeta M_v$. Summing and using \eqref{eq:sum-M} costs at most
$\Lambda/4$. On class (ii), the maximum is bounded in expectation by the
total influence, namely $M_v<M_0$. Since $|\mathcal D|\le2^m$, these
contexts cost at most $\Lambda/4$. On class (iii), the same total-mass
bound is at most $K$ per node.
Equation~\eqref{eq:even-influence} bounds all even contributions by
$\Lambda/4$. Let $N_{\mathrm{iii}}$ count the nodes whose induced
contexts lie in class~(iii). Comparing all four contributions with
\eqref{eq:large-max-sum} gives
\[
 \Lambda
 \le K2^{-s}+\zeta K+2^m M_0+K\E N_{\mathrm{iii}}
 \le\frac{3\Lambda}{4}+K\E N_{\mathrm{iii}}.
\]
Thus $\E N_{\mathrm{iii}}\ge\Lambda/(4K)$.
At each such node \eqref{eq:wq-mq} supplies a label with
$m_q>\gamma M_0/K=\theta$.
Choosing a uniform node from the slice loses at most a factor $2^m$ and
gives \eqref{eq:extraction-event}. This proves the proposition.
\end{proof}
